\section{Variables}

\subsection{Tipo}

La variable independiente es el \textbf{sistema predictivo de precios de exportación de palta Hass}, compuesto por los modelos de precio, margen y riesgo, el simulador de escenarios, el análisis histórico y los reportes de ciclos. La variable dependiente es el \textbf{apoyo a la toma de decisiones de siembra}, evaluado mediante precisión predictiva, comprensión de la información, reducción de incertidumbre y utilidad percibida.

Como variables de entrada se utilizarán fecha, valor FOB, peso neto, destino, temporada, costos agrícolas, superficie y producción. El precio FOB unitario, la productividad, el ingreso, el costo total, el margen y la ganancia serán variables calculadas y deberán conservar trazabilidad respecto de los campos originales.

\begin{table}[H]
\centering
\caption{Variables analíticas y fuente prevista}
\small
\begin{tabularx}{\textwidth}{p{3.3cm}p{3.1cm}X}
\toprule
\textbf{Variable} & \textbf{Campo o cálculo} & \textbf{Fuente y uso} \\
\midrule
Precio FOB unitario & Valor FOB / peso neto & Aduanet; variable objetivo del pronóstico. \\
Volumen exportado & PESO\_NETO & Aduanet; presión de oferta y agregación temporal. \\
Fecha y temporada & FECHA, semana, mes & Aduanet e INEI; orden temporal y estacionalidad. \\
Destino & PAIS\_DESC & Aduanet; segmentación por mercado. \\
Costos agrícolas & P235\_VAL, P237\_VAL, P239, P241 & INEI; estimación del costo productivo. \\
Producción & P219\_CANT\_1, P219\_CANT\_2, P219\_EQUIV\_KG & INEI; productividad y escenarios. \\
Margen y ganancia & Variables calculadas & Simulador; diferencia entre ingresos y costos. \\
\bottomrule
\end{tabularx}
\end{table}

\subsection{Operacionalización}

La matriz de operacionalización se presenta en el Anexo A. Allí se especifican las definiciones conceptual y operacional, dimensiones, indicadores y escalas de medición de las variables independiente y dependiente.
